% Chapter 14 -- FLAIR and InitechDOS (Part II, Reference).
% Sources (Law 1): os/milton/kmain.c finder_exit_dos (the alert, the open
% application asked to quit, DESKTOP.DB committed, the desktop's hooks
% removed, the console handed back, a TRANSIENT COMMAND.COM run as the
% permanent command processor: there is no way back to FLAIR but a restart;
% "TPSREP runs only from this prompt"); os/milton/command.c (the startup
% date and time questions when the disk has no AUTOEXEC.BAT; $P$G; DIR,
% CD, the external programs found in the current directory or by the path
% given); harness/fixture_volumes.mk (the data volume: no AUTOEXEC.BAT, the
% programs in APPS); the InitechDOS books for the commands themselves.
% GATES: test-flair-exit-to-dos test-flair-tps-journey test-shell test-dir test-pubs-flair-screens
\Gates{test-flair-exit-to-dos, test-flair-tps-journey, test-shell, test-dir,
  test-type,
  test-pubs-flair-screens}
\chapter{FLAIR and InitechDOS}
\label{ch:dos}
\index{InitechDOS}

FLAIR runs on InitechDOS.  This chapter tells you how the things you see in
FLAIR appear at the InitechDOS command prompt, and how to use the programs
that run only there.

\section{The same disk, two views}
\index{drive A}\index{directories}
\begin{center}\small
\begin{tabular}{@{}p{46mm}p{53mm}@{}}
\textbf{In FLAIR} & \textbf{At the InitechDOS prompt}\\\hline
the \texttt{INITECH} icon & drive A\\
the \texttt{Drive A Files} window & the root directory, \file{A:\textbackslash}\\
a folder & a directory\\
the \file{APPS} folder & the directory \file{A:\textbackslash APPS}\\
an icon's name & the file name, in capital letters\\
the Trash & the directory \file{A:\textbackslash TRASH}\\
the arrangement of the desktop & the file \file{A:\textbackslash DESKTOP.DB}\\
\end{tabular}
\end{center}
Every file you see in FLAIR is an ordinary InitechDOS file, and the
InitechDOS commands work on it as on any other.  Folders you make with
\menuitem{New Folder} are directories, and files you throw away are kept in
the \file{TRASH} directory until you empty the Trash.

\section{Going to the InitechDOS prompt}
\index{Exit to InitechDOS command}
Choose \menuitem{Exit to InitechDOS} from the Finder's \menuname{Special}
menu and click \texttt{Exit} (Chapter~\ref{ch:starting}).  FLAIR asks the open application
to quit, saves the arrangement of the desktop and hands the screen to
InitechDOS.

Because the data disk has no \file{AUTOEXEC.BAT} file, InitechDOS asks for
the date and the time.  Press \key{Enter} twice to keep them.  The prompt
\texttt{A:\textbackslash>} shows that you are at the top level of drive A.

\section{Looking at the disk from the prompt}
\begin{TaskSteps}
\Step Type \type{CLS} and press \key{Enter}.
\Result{The screen is cleared.}
\Step Type \type{DIR} and press \key{Enter}.
\Result{InitechDOS lists the top level of the data disk: \file{README.TXT},
the directory \file{APPS}, and the directory \file{TRASH} and the file
\file{DESKTOP.DB} that FLAIR does not show in its windows.}
\end{TaskSteps}

\FlairText[fig:dosdir]{u_dosdir}{The top level of the data disk at the
InitechDOS prompt}

\section{Looking at a document from the prompt}
\index{TYPE}
A text document you save with InitechText can be shown at the prompt with
the \file{TYPE} command.
\begin{TaskSteps}
\Step Type \type{TYPE README.TXT} and press \key{Enter}.
\Result{The contents of \file{README.TXT} appear on the screen, followed by
the prompt.}
\end{TaskSteps}

\FlairText[fig:ktype]{k_type}{README.TXT shown with TYPE}

A document longer than the screen scrolls past; to read a long document,
open it in InitechText instead.

\section{Running a program from the APPS folder}
\index{CHKDSK}
The programs in \file{APPS} whose names end in \file{.COM} run at the
prompt.  Change to the directory first, then type the program's name.
\begin{TaskSteps}
\Step Type \type{CLS} and press \key{Enter}.
\Step Type \type{CD APPS} and press \key{Enter}.
\Result{The prompt changes to \texttt{A:\textbackslash APPS>}.}
\Step Type \type{CHKDSK} and press \key{Enter}.
\Result{CHKDSK checks the data disk and reports on it.}
\end{TaskSteps}

\FlairText[fig:doschk]{u_doschk}{CHKDSK run from the APPS directory}

The \doctitle{InitechDOS Reference} describes \file{ATTRIB}, \file{CHKDSK},
\file{EDLIN}, \file{LABEL}, \file{TREE} and \file{XCOPY}.  The TPS Report
Generator, \file{TPSREP}, also runs only from this prompt; its book,
\doctitle{TPS Report Generator Operating Procedures}, tells you how.

\section{Running the TPS Report Generator}
\index{TPS Report Generator}
The TPS Report Generator, \file{TPSREP.COM}, runs only at the InitechDOS
prompt, from the \file{APPS} directory, where its database \file{TPS.DBF}
and its settings file \file{TPSPARM.DAT} are kept.
\begin{TaskSteps}
\Step Leave FLAIR with \menuitem{Exit to InitechDOS}, and press
\key{Enter} twice to keep the date and the time.
\Step Type \type{CD APPS} and press \key{Enter}.
\Step Type \type{TPSREP TPS.DBF} and press \key{Enter}.
\Result{The TPS Report Generator starts and gives its first instruction.}
\end{TaskSteps}

\FlairText[fig:otps]{o_tps}{The TPS Report Generator, started from the APPS
directory}

Follow the procedures in \doctitle{TPS Report Generator Operating
Procedures}, which describes every step of the report and the files it
writes.  When the program ends, the prompt \texttt{A:\textbackslash APPS>}
returns.

\section{Returning to FLAIR}
\index{FLAIR!returning to}
FLAIR does not start again from the prompt.  To return to it, restart the
computer: press its reset button, or turn it off and on.  FLAIR starts with
the arrangement of the desktop it saved when you left.

\begin{MacWarning}
Do not delete \file{DESKTOP.DB} or the \file{TRASH} directory at the
InitechDOS prompt, and do not move files in or out of \file{TRASH} there.
FLAIR makes a new \file{DESKTOP.DB} if the file is missing, but the
arrangement of your desktop is lost.
\end{MacWarning}

\section{Copying files}
\index{copying files}
The Finder moves files but does not copy them.  To make a copy of a file,
use the \file{COPY} command at the prompt; to copy a directory with
everything in it, use \file{XCOPY}.  The \doctitle{InitechDOS User's Guide}
shows you how.
